\begin{textsample}{POS Dim 1 – vem\_ed\_01 – Score 5.00 – vem\_ed\_01/t002.txt}  \label{ex:f1_pos_096}
RESILIÊNCIA INTERNACIONAL Projetos Colaborativos Internacionais (PCIs) das Fatecs superam as adversidades do isolamento social e contribuem para formar tecnólogos globais EXPERIÊNCIA E DETERMINAÇÃO Um dos PCIs mais longevos ocorre entre a Fatec Americana e a State University of New York (Suny Ulster).
Conduzido desde 2013 por Carlos Augusto Amaral Moreira, tem como objetivo simular o lançamento de um \textbf{produto} ou serviço brasileiro nos EUA (e vice-versa) e já ganhou prêmio do Guia do Estudante / Santander."Há dificuldades, mas sempre buscamos superar", afirma Moreira.
Geralmente as equipes mistas reúnem 6 brasileiros (calouros de Gestão Empresarial, matriculados na \textbf{disciplina} Administração Geral) e 4 dos EUA (estudantes de Principles of Management), com um líder de \textbf{cada} nacionalidade em \textbf{cada} grupo.
A \textbf{partir} do exame de proficiência, os calouros \textbf{que} têm melhor \textbf{nível} de proficiência em inglês se distribuem, para ajudar os colegas e assumirem as videoconferências.
Os demais privilegiam texto escrito.
Depois das \textbf{interações} iniciais, \textbf{cada} grupo aborda um dos seis temas: cultura, economia, demografia, \textbf{política}, história, tendências da juventude.
Na \textbf{terceira} fase do projeto, \textbf{que} dura de 8 a 10 semanas, 3 equipes desenvolvem \textbf{produto} para os EUA e outras 3 para o Brasil, incluindo uma rápida pesquisa de \textbf{mercado}.
A última etapa envolve planejamento estratégico.
Quando se \textbf{iniciou} o recesso nas Fatecs (23 de março), como \textbf{todo} o corpo docente, parou as atividades, mas combinou com alunos e o professor parceiro de voltarem às atividades assim \textbf{que} \textbf{possível}.
Nos EUA, as aulas passaram para a modalidade remota.
No retorno às aulas por videoconferência nas Fatecs (22 de abril), Moreira ajustou o calendário, eliminando pesquisa de \textbf{mercado} e \textbf{reduzindo} o planejamento estratégico para 10 dias.
Deu certo: os seis grupos entregaram as apresentações.
Esse projeto \textbf{foi} um dos dois únicos \textbf{que} permaneceu \textbf{ativo}, mesmo com a interrupção do calendário acadêmico.
O outro, também tocado por Moreira, só \textbf{que} com a Florida State University (FSU) envolve mapeamento de \textbf{diferenças} culturais entre Brasil e EUA.
Na FSU, participam os matriculados em Intercultural Competence. \textbf{São} alunos de Química, Engenharia, Biologia, entre outros cursos; na Fatec Americana, \textbf{são} estudantes de Gestão Empresarial, na \textbf{disciplina} Comportamento Organizacional.
Os grupos mistos envolvem 4 \textbf{pessoas}, \textbf{que} falam semanal ou quinzenalmente por chamada de vídeo no WhatsApp.
Os primeiros \textbf{diálogos} \textbf{começaram} em fevereiro e até meados de março já haviam elaborado um relatório sobre dimensões culturais, conforme noções de Geert Hofstede (psicólogo holandês falecido em fevereiro deste ano).
As atividades \textbf{foram} retomadas em 22 de abril, com a realização de um vídeo de 4 minutos \textbf{comparando} \textbf{diferenças} e semelhanças entre os dois países."No primeiro dia do retorno, 80 \% dos alunos participaram da videoconferência, \textbf{foi} bacana ver \textbf{que} os alunos topam o desafio".

% matched lemmas: ativo, cada, começar, comparar, diferença, disciplina, diálogo, iniciar, interação, mercado, nível, partir, pessoa, político, possível, produto, que, reduzir, ser, terceiro, todo
\end{textsample}
